\section{Mechanism: Why the Temporal Term Shifts Ranking}

\label{sec:mechanism}

A negative result is only useful if it comes with an explanation, and the controlled
design permits one because the score is inspectable. This section explains why the
deployed hybrid loses to a lexical baseline, using the recorded score decomposition
rather than a post-hoc narrative. The analysis describes the existing scoring function;
it does not propose a tuned replacement, and the illustrative case below follows a
disclosed deterministic selection rule rather than being chosen to support the
argument.

\subsection{The pairwise condition}

\label{sec:mech_condition}

In the replay protocol the memory timestamps are the original conversation timestamps
and the reference clock is the question date. Frequency and importance are fixed in
the controlled comparison, so pairwise order under the deployed hybrid is determined by

\begin{equation}
  S_i = 0.4\,J(q,m_i) + 0.3\,D_i + C,
  \label{eq:mech_score}
\end{equation}

where $J$ is lexical Jaccard similarity, $D$ is the implemented temporal decay, and $C$
collects the constant frequency and importance contributions. A relevant older record
outranks a fresher record precisely when

\begin{equation}
  0.4\,(J_{\mathrm{old}} - J_{\mathrm{fresh}})
  \;>\; 0.3\,(D_{\mathrm{fresh}} - D_{\mathrm{old}}).
  \label{eq:mech_ineq}
\end{equation}

Equation~\ref{eq:mech_ineq} is the whole mechanism in one line. The left side is the
lexical advantage of the older, more relevant record; the right side is the freshness
advantage of the newer, less relevant record. Because the temporal term is
query-independent, the right side does not depend on what is being asked. The deployed
weights make the comparison acute: with $w_s=0.4$ and $w_t=0.3$, a record must be more
than roughly three-quarters again as lexically similar as its fresher competitor before
its relevance advantage can be expected to dominate a full-range freshness gap. When
decay is close to saturated, $D_{\mathrm{fresh}}-D_{\mathrm{old}}$ can approach $0.3$
in weighted units while the lexical gap is bounded by $0.4$ in the best case and is
typically much smaller, since Jaccard overlap between a short question and a long
segment is low for every candidate.

\subsection{Consequences for the bank}

\label{sec:mech_consequences}

Three consequences follow, and all three are visible in the recorded data. First, the
deficit should be largest where the query is lexically specific and the relevant fact is
old, which is the single-session-user stratum; that is where the measured gap is widest,
with BM25 at 100.0\% and the hybrid at 11.1\%. Second, removing only the temporal term
should recover most of the loss while leaving everything else intact, and it does:
hybrid minus hybrid\_no\_time is $-24.2$, $-27.0$ and $-33.3$ points at the three
budgets, with Holm-adjusted $p\le 0.0022$ throughout. Third, the evidence-hit rate
should fall with the accuracy rate, because the items displaced from the prompt are
exactly the relevant older ones; that is what the 18.1\% against 75.0\% comparison at
512 tokens shows.

The mechanism also explains the shape of the budget trend. One might expect extra
budget to rescue the hybrid, but it does not, because the temporal term reorders the
candidate list at every budget rather than merely truncating it. Adding tokens lets the
hybrid include more items, and the items it prefers to include are still the fresher
ones. The absolute accuracy therefore stays near 33--37\% at all three budgets while
BM25 stays near 66--71\%.

\FloatBarrier

\subsection{A recorded illustrative case}

\label{sec:mech_case}

The recorded cases use a deterministic example-selection rule and do not add
independent evidence to the paired tests; they exist to show the mechanism operating on
real records rather than on a constructed example. The pattern they illustrate is one in
which a question's gold evidence lies in an older session while a newer session shares
surface vocabulary with the question, so lexical similarity and recency point at
different records. Under Equation~\ref{eq:mech_ineq} the newer record wins the
comparison even though it does not contain the answer, the delivered context omits the
gold evidence, and the resulting answer is incorrect. Removing the temporal term, or
replacing Jaccard with BM25, changes which record is selected and which evidence reaches
the prompt, which is why the same candidate bank yields different accuracy under
different ordering rules. The appendix lists the recorded cases with their session
identifiers so that the selection can be checked against the raw records.

It is worth being explicit about what this explanation does not establish. Equation
~\ref{eq:mech_ineq} is an account of the existing score, not a claim that the weights
were badly chosen in some objective sense; the weights are the implementation's
defaults and were not fitted on the evaluation questions. The replay uses original source
timestamps and the question date rather than assigning a common ingestion timestamp, and
bulk ingestion that stamps every record with the same time is a different scenario that
this experiment does not evaluate. Frequency and importance are held fixed, so the
analysis isolates the interaction of lexical relevance and temporal decay rather than the
full four-term score.

\FloatBarrier
